\section{literature review}

Pick and place requirement is one of the common applications of robotics in modern industries. Research has demonstrated robotic arms which are capable of detecting, grasping and transferring the objects between two different locations. One of the examples is an intelligent six-axis robotic arm which used for pick-and-place operations which is used for object recognition and demonstrated the integration of object detection with robotic manipulation.\cite{tung2023}

The kinematics are the main concept of the robotic arm control; it ensures the movements of the end effector and according to that the configuration of the individual joints in the system. Inverse kinematics is used to determine the joints’ angles and other variables to position the end effector at the target location accurately. This concepts investigates both analytical and learning based approaches and analytical methods providing solutions based on the mathematical model of the robotic manipulator.\cite{wagaa2023}

Automated pick-and-place systems consist of various units and object detection is one of the important ones. It should identify the object correctly and determine when and where the object is available for manipulation. Previous research has investigated camera-based object detection and machine-learning approaches. Specially vision-based systems are developed for this purpose in most modern robotic systems and provide the coordinates to the robotic arm for the pick-and-place operations. \cite{kumar2014}

Most robotic designs are integration of object detection, kinematic computation, actuator control within embedded platforms etc. A recent six DOF robotic arm study which integrated object detection and robot control on a resource constrained microcontroller is one of the examples for the demonstration of the possibility of performing perception and manipulation using an embedded architecture.\cite{almaliki2026}
